% !TeX root = ../../Book2.tex
\section{限制算子}
\label{b2:sec:2301}

如果 \(f\in L^1(\mathbb R^d)\)，其 Fourier 变换是一个确定的连续函数，可以在任何曲面上取值。若只知道 \(f\in L^2\)，Plancherel 给出的却是相对于 \(d\) 维 Lebesgue 测度的等价类。一个光滑超曲面的 \(d\) 维测度为零，在这个集合上任意改值都不改变等价类。因此，限制问题首先是一个定义问题：是否存在一个连续算子，与光滑输入的曲面取值一致？

\ProseParagraphBreak
这个问题与通常的迹问题相似，但所控制的范数不同。频率变量的 Sobolev 正则性可以产生迹；限制定理则希望只用原函数的普通 \(L^p\) 范数，并由曲面的振荡几何取得额外控制。两种办法的条件和作用不能混用。

\subsection{\texorpdfstring{\(\widehat f|_S\)}{Fourier 变换在 S 上的限制} 的意义}
\label{b2:subsec:230101}

先设 \(d\geq2\)，\(S\subset\mathbb R^d\) 是非空紧致光滑嵌入超曲面，取面积测度 \(\sigma\)。仍用
\(\widehat f(\xi)=\int f(x)e^{-2\pi\mathrm i x\cdot\xi}\,dx\)。
对 \(f\in\mathcal S(\mathbb R^d)\)，定义
\[
 Rf\coloneqq\widehat f|_S.
\]
\symidx{\(Rf\)：曲面 Fourier 限制}
一个 \(L^p\) 到 \(L^r\) 的\term[fourier-xianzhi]{Fourier 限制估计}[Fourier restriction estimate]是指
\[
 \|Rf\|_{L^r(S,d\sigma)}\leq C\|f\|_{L^p(\mathbb R^d)},
 \quad f\in\mathcal S(\mathbb R^d).
\]
这里 \(1\leq p<\infty\)、\(1\leq r\leq\infty\)，常数与 \(f\) 无关。对于带紧支撑光滑密度的局部曲面片，也使用同一定义，只需把面积测度换成指定的有限测度。

\ProseParagraphBreak
为什么这个估计足以解决定义问题？若 \(f_n\in\mathcal S\) 在 \(L^p\) 中趋于 \(f\)，则
\[
 \|Rf_n-Rf_m\|_r\leq C\|f_n-f_m\|_p\longrightarrow0.
\]
目标 \(L^r(S)\) 是 Banach 空间，包括 \(r=\infty\) 的情形，故有极限。用同一不等式比较两组逼近，极限不依赖选择，从而定义了唯一的有界延拓 \(R:L^p\to L^r(S)\)。这个 \(Rf\) 是曲面测度意义下的等价类；它并不要求先给环境 Fourier 等价类选一个逐点代表。

\ProseParagraphBreak
最简单的估计来自绝对积分：
\[
 |\widehat f(\xi)|\leq\|f\|_1,\quad
 \|Rf\|_r\leq\sigma(S)^{1/r}\|f\|_1,
\]
其中 \(1/\infty=0\)。所以 \(p=1\) 对所有目标指数都成立。一般 \(p\) 的延拓若已建立，则对 \(f\in L^1\cap L^p\) 仍与连续 Fourier 变换的取值一致：选同时在两种范数中逼近的 Schwartz 函数，\(L^1\) 收敛使 Fourier 变换在整个频率空间一致收敛，而限制估计给出曲面范数收敛，两种极限必相同。

\ProseParagraphBreak
相反，\(p=2\) 的困难不能仅靠选一个较弱的曲面范数消除。下面把薄壳反例写成实际的 Schwartz 函数，避免仅从零测集问题作形式判断。

\Needspace{6\baselineskip}
\begin{proposition}[薄壳输入]\label{b2:prop:23-thin-shell}
设 \(d\geq2\)，\(S\subset\mathbb R^d\) 是非空光滑嵌入超曲面，取面积测度。对任意 \(1\leq r\leq\infty\)，不存在常数 \(C\)，使
\[
 \|\widehat f|_S\|_{L^r(S)}\leq C\|f\|_2
 \quad\text{对全部 }f\in\mathcal S(\mathbb R^d)\text{ 成立}.
\]
即使只要求控制 \(S\) 上一个非空的小坐标片，也没有这样的界。
\end{proposition}
\begin{proof}
经平移和正交变换，可把一小片写成
\(S=\{(\eta,h(\eta)):\eta\in U\}\)，其中 \(U\subset\mathbb R^{d-1}\) 开，\(h\) 光滑。在这一局部表示中选非零
\(b\in C_c^\infty(U)\)，再取 \(\chi\in C_c^\infty(\mathbb R)\)，使 \(\chi(0)=1\)。对充分小的 \(\epsilon>0\)，令
\[
 \widehat f_\epsilon(\eta,\xi_d)
 =b(\eta)\chi\bigl((\xi_d-h(\eta))/\epsilon\bigr),
\]
并在该坐标邻域外延零。支撑可选在邻域内部，故这确实是一个全空间的光滑紧支撑函数，\(f_\epsilon\in\mathcal S\)。
在曲面片上的值恒为 \(b(\eta)\)，因此其 \(L^r\) 范数是一个不随 \(\epsilon\) 变化的正数。

另一方面，作 \(\xi_d=h(\eta)+\epsilon t\) 换元，Jacobi 因子为 \(\epsilon\)，Plancherel 给出
\[
 \|f_\epsilon\|_2^2
 =\|\widehat f_\epsilon\|_2^2
 =\epsilon\|b\|_{L^2(\mathbb R^{d-1})}^2
       \cdot\|\chi\|_{L^2(\mathbb R)}^2.
\]
假定的限制界会把固定正数控制为 \(O(\epsilon^{1/2})\)，令 \(\epsilon\downarrow0\) 即矛盾。
\end{proof}

薄壳越来越窄，环境平方质量趋零，曲面上的值却没有变小。若增加空间矩的控制，频率函数便获得法向正则性，情况就不同了。下面的球面模型把两种理论的差别显示出来。

\Needspace{6\baselineskip}
\begin{proposition}[带空间矩的平方迹]\label{b2:prop:23-moment-trace}
设 \(d\geq2\)，单位球面 \(S^{d-1}\) 取面积测度。对每个
\(f\in\mathcal S(\mathbb R^d)\)，有
\[
 \|\widehat f|_{S^{d-1}}\|_2^2
 \leq C_d\|f\|_2\cdot
       \bigl(\|f\|_2+\||x|f\|_2\bigr).
\]
因此球面限制唯一连续延拓到
\[
 X\coloneqq\{f\in L^2(\mathbb R^d):|x|f\in L^2(\mathbb R^d)\},
 \quad \|f\|_X=\|f\|_2+\||x|f\|_2,
\]
并保持上述估计。
\end{proposition}
\begin{proof}
对固定 \(\omega\in S^{d-1}\)，令 \(F(r)=\widehat f(r\omega)\)。
对任意 \(s\in(1/2,3/2)\)，微积分基本定理给出
\[
 |F(1)|^2\leq |F(s)|^2+
      2\int_{1/2}^{3/2}|F(r)|\cdot|F'(r)|\,dr.
\]
把 \(s\) 在长度为一的区间上积分，再对 \(\omega\) 积分。因为
\(r^{d-1}\) 在这个区间有正的统一上下界，并且
\(|F'(r)|\leq|\nabla\widehat f(r\omega)|\)，在壳
\(A=\{\xi:1/2<|\xi|<3/2\}\) 上得到
\[
 \|\widehat f|_{S^{d-1}}\|_2^2
 \leq C_d\|\widehat f\|_{L^2(A)}^2
       +C_d\|\widehat f\|_{L^2(A)}
                \cdot\|\nabla\widehat f\|_{L^2(A)}.
\]
Plancherel 和
\(\partial_{\xi_j}\widehat f=-2\pi\mathrm i\,\widehat{x_jf}\)
给出
\(\|\nabla\widehat f\|_2=2\pi\||x|f\|_2\)，即得估计。

最后，\(C_c^\infty\) 在 \(X\) 中稠密。先截断 \(f\) 的空间支撑，由两项平方可积性控制尾部；对紧支撑输入再卷积光滑近似恒等核。此时
\[
 x_j(\rho_\epsilon*f)
 =\rho_\epsilon*(x_jf)+(x_j\rho_\epsilon)*f,
\]
后一个核的 \(L^1\) 范数为 \(O(\epsilon)\)，故两项加权范数都收敛。估计应用于逼近差，使球面迹为 Cauchy 列，完成延拓。
\end{proof}

带矩范数提供的是 \(\widehat f\) 在频率变量中的一阶平方正则性。
Stein–Tomas 定理将改用一个小于 \(2\) 的普通 Lebesgue 指数，不要求空间矩。
这两种输入条件并不相互等同：在无限测度空间上，单纯改变 \(p\) 不能直接替代位置权重。

\subsection{曲面测度}
\label{b2:subsec:230102}

曲面限制必须指定相对于哪一个测度计算范数。在图形坐标
\(\xi=(\eta,h(\eta))\) 中，面积测度为
\[
 d\sigma(\xi)=\sqrt{1+|\nabla h(\eta)|^2}\,d\eta.
\]
这个因子是切向面积的 Jacobi 因子。若改用
\(d\mu=a\,d\sigma\)，其中 \(a\geq0\) 光滑且紧支撑，则在 \(a\) 有正上下界的紧片上，两个有限指数范数等价。密度若消失，反例的集中质量可能随之变小，不能仍按恒定正密度计算必要条件。以后局部化曲面时，都把密度和面积因子写进振幅中。

\ProseParagraphBreak
球面的极坐标公式具体说明了面测度怎样由薄壳得到：
\[
 \int_{\mathbb R^d}F(\xi)\,d\xi
 =\int_0^\infty r^{d-1}
       \int_{S^{d-1}}F(r\omega)\,d\sigma(\omega)\,dr.
\]
例如对连续紧支撑的 \(F\)，有
\[
 \lim_{\epsilon\downarrow0}
 \frac1{2\epsilon}\int_{1-\epsilon<|\xi|<1+\epsilon}F(\xi)\,d\xi
 =\int_{S^{d-1}}F(\omega)\,d\sigma(\omega).
\]
证明只需在固定紧壳上用一致连续性，并注意
\(r^{d-1}\to1\) 一致成立。这里已经除以厚度 \(2\epsilon\)；
没有这个归一化，壳积分本身趋于零。

\ProseParagraphBreak
对于一般定义函数 \(\rho\)，假设 \(\nabla\rho\neq0\) 于
\(S=\{\rho=0\}\) 附近，沿法向换元得到分布记号
\[
 \int F(\xi)\delta(\rho(\xi))\,d\xi
 =\int_S\frac{F(\xi)}{|\nabla\rho(\xi)|}\,d\sigma(\xi).
\]
它可先在坐标片内将 \(\rho\) 作为一个坐标验证，再用有限单位分解拼接。
特别地，
\(\delta(1-|\xi|^2)=\sigma/2\) 在单位球面上成立。
后面解析族出现的二分之一来自此处的换元，不来自 Fourier 归一化。
上述公式对测试函数成立，并不自动控制任意 \(L^p\) 输入的薄壳极限；取得这种控制仍需要限制估计。

\subsection{曲率}
\label{b2:subsec:230103}

对图形曲面 \((\eta,h(\eta))\)，选定单位法向后，其 Gaussian 曲率为
\[
 K(\eta)=\frac{\det D^2h(\eta)}
                 {(1+|\nabla h(\eta)|^2)^{(d+1)/2}},
\]
符号可能随法向约定改变，但非零性不变。因此本章使用的条件是
\(\det D^2h\neq0\)，并不要求所有主曲率同号。
在一个紧支撑片上，行列式绝对值有正下界；这使\chapref{b2:ch:21}的非退化驻相估计具有统一常数。
球面在每个点都满足这个条件。

\ProseParagraphBreak
曲率的作用有两个方面。对于光滑密度，\(d-1\) 个切向变量的二次振荡带来
\((1+|x|)^{-(d-1)/2}\) 衰减。对于小帽数据，曲面离开切平面的高度却只有帽半径的平方，这允许输出在法向保持更长距离的相干。前者用于证明充分估计，后者产生必要条件；两者是同一二次几何的不同后果。

\ProseParagraphBreak
平面上的局部问题已经能显示曲率消失的影响。取
\[
 S_0=\{(\eta,0):|\eta|<1\}\subset\mathbb R^d
\]
及其面积测度。选 \(g\in\mathcal S(\mathbb R^{d-1})\)，使
\(\widehat g\) 在该片上不恒为零，另选 \(h\in\mathcal S(\mathbb R)\)，使
\(\int h\neq0\)。令
\[
 f_L(x',x_d)=g(x')h(x_d/L),\quad L\geq1.
\]
对 \(1<p<\infty\)，
\[
 \|f_L\|_p=L^{1/p}\|g\|_p\cdot\|h\|_p,\quad
 \widehat f_L(\eta,0)=L\widehat g(\eta)\int h.
\]
任何固定 \(1\leq r\leq\infty\) 的限制范数都按 \(L\) 增长，比输入的
\(L^{1/p}\) 快。故这个有界平面片也没有 \(p>1\) 的限制界。
无穷大的平面面积不是反例的原因；法向上没有振荡损失才是关键。

\subsection{缩放}
\label{b2:subsec:230104}

单位球面不是尺度不变对象。把它放大时，既改变面测度，又改变 Fourier 变换中的空间尺度，两个因子都要保留。

\Needspace{6\baselineskip}
\begin{proposition}[曲面半径的缩放]\label{b2:prop:23-surface-scaling}
设 \(d\geq2\)，\(S\) 是紧致光滑超曲面，\(1\leq p<\infty\)、
\(1\leq r\leq\infty\)。假设相对于面积测度的限制算子
\(L^p(\mathbb R^d)\to L^r(S)\) 的最佳常数为 \(C\)。对 \(R>0\)，
\(S_R=RS\) 也取面积测度，则其最佳常数恰为
\[
 C_R=C R^{(d-1)/r-d/p'},
 \quad \frac1p+\frac1{p'}=1.
\]
\end{proposition}
\begin{proof}
取 \(g(x)=f(x/R)\)。直接换元给出
\[
 \widehat g(\xi)=R^d\widehat f(R\xi),
 \quad \|g\|_p=R^{d/p}\|f\|_p.
\]
同时 \(d\sigma_R(R\xi)=R^{d-1}d\sigma(\xi)\)，所以
\[
 \|\widehat f\|_{L^r(S_R)}
 =R^{(d-1)/r-d}\|\widehat g\|_{L^r(S)}
 \leq C R^{(d-1)/r-d/p'}\|f\|_p.
\]
\(r=\infty\) 时第一因子按 \(R^0\) 理解。缩放是 Schwartz 类的双射，
反向代入近乎达到最佳常数的输入，便知这个因子不只是一个上界。
\end{proof}

例如三维球面 Stein–Tomas 端点的 \(p=4/3,r=2\) 将产生
\(R^{1/4}\)，而不是与半径无关的常数。单位球面估计一旦成立，缩放便确定了其他半径所需的因子。若把面测度改为总质量一的概率测度，还须再乘
\(R^{-(d-1)/r}\)；使用不同测度归一化的公式时，应先对齐这一约定。
